% Source-derived chapter 10: Applications for rational points
% Original source: chabauty-coleman-bound-surfaces
\chapter{Applications for rational points}
\label{chabauty-coleman-bound-surfaces-chapter-10}
\source{chabauty-coleman-bound-surfaces}

\begin{remark}[Source section]
\label{chabauty-coleman-bound-surfaces-chapter-10-source}
\source{chabauty-coleman-bound-surfaces}
\source{chabauty-coleman-bound-surfaces}
This chapter reproduces the corresponding section of the retrieved
LaTeX source, including its definitions, statements, examples, and
proofs. It has not yet been translated into Lean.
\notready
\end{remark}

% The source section title was: Applications for rational points
\label{SecApplications}
\source{chabauty-coleman-bound-surfaces}


%%%%%%%%%
%%%%%%%%%
%%%%%%%%%

\subsection{Rational points on surfaces} \label{SecApp1} We begin with the following generalization of Theorems \ref{ThmIntroA} and \ref{ThmIntroB} to the case of an arbitrary number field.
\source{chabauty-coleman-bound-surfaces}


%%
%%
\begin{theorem}
\source{chabauty-coleman-bound-surfaces}
\notready\label{ThmFinalMain} Let $F$ be a number field, let $p$ be a prime number, let $\pfrak$ be a prime of $F$ above $p$, let $\efrak$ be the ramification index of $\pfrak$, let $\kk$ be the residue field of $\pfrak$, and let $q=\#\kk$. We assume that 
\source{chabauty-coleman-bound-surfaces}
$$
p>\max\{\efrak+1, \exp(\efrak/\exp(1))\}.
$$
Let $X$ be a smooth projective surface contained in an abelian variety $A$ of dimension $n\ge 3$, both defined over $F$ and having good reduction at $\pfrak$. Let $X'$ and $A'$ be the corresponding reductions modulo $\pfrak$. Suppose  $\rk A(F)\le 1$ and that either of the following conditions holds:
\begin{itemize}
\item[(i)] $n=3$, $X$ is of general type, $X'$ contains no elliptic curves over $\kk^{alg}$, and 
$$
p> (128/9)\cdot c_1^2(X)^2.
$$ 
\item[(ii)] $A'$ is simple over $\kk^{alg}$ and there is an ample divisor $H$ on $A$ such that
$$
p> \max\left\{3 c_1^2(X)+2,\frac{n!\cdot (3\deg_H(X) + \cdeg_H(X))^n}{n^n\cdot \deg(H^n)}\right\}.
$$
\end{itemize}
Then $X(F)$ is finite and we have
$$
\# X(F) \le \# X'(\kk) + \left(1-\frac{\efrak}{p-1}\right)^{-1}\cdot (q+4q^{1/2}+3)\cdot c_1^2(X).
$$
\end{theorem}
\begin{proof} By Theorem \ref{ThmMain}, after base-change to the completion $K=F_\pfrak$ and taking $G=A(F)$. 
\end{proof}
%%
%%



%%%%%%%%%
%%%%%%%%%
%%%%%%%%%

\subsection{Primes of geometrically simple reduction} \label{SecApp2} The next result is due to Chavdarov \cite{Chavdarov}.
\source{chabauty-coleman-bound-surfaces}
%%
%%
\begin{lemma}
\source{chabauty-coleman-bound-surfaces}
\notready\label{LemmaCha} Let $A$ be an abelian variety over $\Q$ of odd dimension $n$ satisfying $\End (A_\C)=\Z$. There is a set of primes $\Pcal$ of density $1$ in the primes such that for every $p\in \Pcal$ the reduction of $A$ modulo $p$ is geometrically simple.
\source{chabauty-coleman-bound-surfaces}
\end{lemma}
%%
%%
With this at hand, we have:

\begin{proof}[Proof of Corollary \ref{CoroDim3}] Let $\Pcal$ be the set of primes afforded by Lemma \ref{LemmaCha} applied to $A$ (as $\dim A=3$ is odd), discarding the primes of bad reduction for $X$. For every $p\in \Pcal$ we have that $X\otimes \F_p^{alg}$ does not contain elliptic curves since $A\otimes \F_p^{alg}$ is simple. We conclude by Theorem \ref{ThmIntroA}.
\end{proof}

From the argument it is clear that we can get a similar result over number fields and abelian varieties of dimension $n\ge 3$ using Theorem \ref{ThmFinalMain}, provided that a suitable extension of Chavdarov's theorem is applicable to $A$. The general problem regarding abundance of primes of geometrically simple reduction is a conjecture of Murty and Patankar \cite{MurtyPatankar}. See \cite{Zywina} and the references therein for the formulation of this conjecture and available results.

%%%%%%%%%
%%%%%%%%%
%%%%%%%%%

\subsection{Quadratic points in curves} \label{SecApp3} We write $\equiv$ for numerical equivalence.
\source{chabauty-coleman-bound-surfaces}

Let $C$ be a smooth,  irreducible, projective curve of genus $g\ge 2$ with Jacobian $J$, over an algebraically closed field $k$ of characteristic different from $2$. The symmetric square of $C$,  denoted by $C^{(2)}$, is a smooth, irreducible, projective surface. The quotient map $\pi:C\times C\to C^{(2)}$ is finite of degree $2$ ramified along the diagonal $\Delta\subseteq C\times C$. Let $D=\pi(\Delta)$ and for $\xi \in C(F^{alg})$ we let $V_\xi=\pi(\{\xi\}\times C)$. As we vary $\xi$, all the divisors $V_\xi$ are numerically equivalent to each other and this numerical equivalence class  is denoted by $V$.  Using adjunction for $\pi$ (see \cite{Hartshorne} Proposition II.8.20) and intersections on $C\times C$ we find


\begin{lemma}
\source{chabauty-coleman-bound-surfaces}
\notready\label{LemmaKSymm2} We have $K_{C^{(2)}}\equiv (2g-2)V-D/2$. Moreover, $V^2=1$, $(D.V)=2$, $D^2=4-4g$, and $c_1^2(C^{(2)})=(4g-9)(g-1)$.
\source{chabauty-coleman-bound-surfaces}
\end{lemma}


Let $d_0$ be degree $2$ divisor on $C$. The map $C\times C\to J$ given by $(x_1,x_2)\mapsto [x_1+x_2-b_0]$ induces a  map $j:C^{(2)}\to J$ birational onto its image (since $g\ge 2$). Let $\Theta$ be a Theta divisor in $J$ and recall that it is ample.  We define $\theta=j^*\Theta$ on $C^{(2)}$. By Lemma 7 in \cite{KouvidakisTAMS93} and Lemma \ref{LemmaKSymm2} we have
 
 \begin{lemma}
\source{chabauty-coleman-bound-surfaces}
\notready\label{LemmaThetaC2} $D\equiv 2((g+1)V-\theta)$. Hence, $\theta\equiv (g+1)V-D/2$ and $K_{\overline{C}^{(2)}}\equiv \theta + (g-3)V$.
\source{chabauty-coleman-bound-surfaces}
\end{lemma}

\begin{corollary}
\source{chabauty-coleman-bound-surfaces}
\notready \label{LemmaSymm2GT} For every curve $C$ of genus $g\ge 3$  we have that $C^{(2)}$ has ample canonical divisor. In particular, $C^{(2)}$ is a surface of general type.
\source{chabauty-coleman-bound-surfaces}
\end{corollary}
\begin{proof} Since $\theta$ is ample and all the $V_{\xi}$ are effective, we get the result from Lemma \ref{LemmaThetaC2}.
\end{proof}


By Lemmas \ref{LemmaKSymm2} and \ref{LemmaThetaC2} (namely, $\theta\equiv (g+1)V-D/2$) we find

\begin{lemma}
\source{chabauty-coleman-bound-surfaces}
\notready\label{LemmaThetaK} $(\theta.K_{C^{(2)}})=2g(g-2)$. 
\source{chabauty-coleman-bound-surfaces}
\end{lemma}


Let $W_2=j(C^{(2)})$. It is clear that

\begin{lemma}
\source{chabauty-coleman-bound-surfaces}
\notready\label{LemmanonHypEll} If $C$ is not hyperelliptic, then $j$ is injective, in which case $C^{(2)}\simeq W_2$.
\source{chabauty-coleman-bound-surfaces}
\end{lemma}

Now consider a field $F$ of characteristic different from $2$, not necessarily algebraically closed and suppose that $C$ is defined over $F$. If the divisor $d_0$ is defined over $F$ then so is the map $j:C^{(2)}\to J$ and its image $W_2$. However, $\Theta$ is only defined in some algebraic extension of $F$ in general. Nevertheless, in this setting we have

\begin{lemma}
\source{chabauty-coleman-bound-surfaces}
\notready\label{LemmaTheta} Assuming the existence of an effective degree $2$ divisor $d_0$ on $C$ defined over $F$, there is a divisor $H$ on $J$ defined over $F$ with $H\equiv 2\Theta$. In particular, $H$ is ample.
\source{chabauty-coleman-bound-surfaces}
\end{lemma}
\begin{proof} Write $d_0=x_1+x_2$ where $x_1,x_2$ are either $F$-rational or of degree $2$ over $F$ in which case they are Galois conjugate. Let $\Theta_i$ be the $(g-1)$-fold pointwise sum of $j(V_{x_i})\subseteq J$ with itself, for $i=1,2$. Then the divisor $H=\Theta_1+\Theta_2$ has the required properties.  
\end{proof}


%%
%%
\begin{proof}[Proof of Corollary \ref{CoroQuadratic}] We keep the previous notation in the case $F=\Q$. We can assume the existence of the effective degree $2$ divisor $d_0$ on $C$ defined over $\Q$, for otherwise $C^{(2)}(\Q)$ is empty. Let us apply Theorem \ref{ThmIntroB} to $A=J$, $n=g$,  and $X=W_2\subseteq J$, which is a smooth projective surface satisfying $W_2\simeq C^{(2)}$ (Lemma \ref{LemmanonHypEll}). Furthermore, the reduction of $C^{(2)}$ modulo $p$ is $(C')^{(2)}$ and the isomorphism $C^{(2)}\simeq X$ induced by $j$ reduces to an isomorphism $(C')^{(2)}\simeq X'$ with $X'$ the reduction of $X$, because $C'\otimes \F_p^{alg}$ is not hyperelliptic. The divisor $H$ is taken as in Lemma \ref{LemmaTheta}. 

By the Poincar\'e formulas (see for instance \cite{KouvidakisTAMS93}) we have $\deg(H^g)=2^g\deg(\Theta^g)=2^g\cdot g!$ and $\deg(H^2.X)=4\deg(\Theta^2.W_2)=4g(g-1)$. By Lemma \ref{LemmaThetaK} we obtain $\deg(H.K_X)=2\deg(\Theta.K_{W_2})=2 (\theta.K_{C^{(2)}})=4g(g-2)$. Therefore,
$$
\frac{g!\cdot (3\deg(H^{2}.X) + \deg(H.K_X))^g}{g^g\cdot \deg(H^g)} = \frac{g!\cdot (4g(4g-5))^g}{g^g\cdot 2^g\cdot g!}=(8g-10)^g.
$$
In view of Lemma \ref{LemmaKSymm2} and the fact that $3c_1^2(X)+2=3(4g-9)(g-1)+2<(8g-10)^g$ for $g\ge 3$, the result follows from Theorem \ref{ThmIntroB}.
\end{proof}

\begin{proof}[Proof of Corollary \ref{CoroQuadratic3}] Similar to the previous argument, using Theorem \ref{ThmIntroA} instead of Theorem \ref{ThmIntroB}. Here, $X=W_2\simeq C^{(2)}$ is of general type by Lemma \ref{LemmaSymm2GT}. Lemma \ref{LemmaKSymm2} and the fact that $g=3$ give $c_1^2(X)=c_1^2(C^{(2)})=6$. Thus, $(128/9)c_1^2(X)^2= 512$ and the least prime $p>512$ is $521$.
\end{proof}

%%%%%%%%%%%%%%%%%%%%%%%%%%%%%%%%%%%%%%
%%%%%%%%%%%%%%%%%%%%%%%%%%%%%%%%%%%%%%
%%%%%%%%%%%%%%%%%%%%%%%%%%%%%%%%%%%%%%
%%%%%%%%%%%%%%%%%%%%%%%%%%%%%%%%%%%%%%
%%%%%%%%%%%%%%%%%%%%%%%%%%%%%%%%%%%%%%
%%%%%%%%%%%%%%%%%%%%%%%%%%%%%%%%%%%%%%
